% !PRINTED 451
\weberpage{451}
%<<CONT
der zweiten Gleichung reduziert sich nur dann auf $D$, wenn sie imprimitiv ist und
\weq{8}{\begin{gathered}
B' - 2A'c \ \text{durch} \ p, \\
A'^{2}c - B'c + C' \ \text{durch} \ p^2 % [sic]
\end{gathered}}
teilbar ist, und dann genügt die entsprechende Zahl (6) gleichzeitig der Gleichung (5) und der Gleichung (2). Dies trifft aber, wie wir gleich zeigen, nur für einen Wert $c$ zu, und folglich läßt sich der betreffende Wert $u$ \gesperrt{rational} durch $v$ ausdrücken.

Denn ist zunächst $p$ ungerade, so hat die erste Kongruenz (8):
\[
2A'c - B' \equiv 0 \pmod{p}
\]
nur \gesperrt{eine} Wurzel, und für diese ist
\[
(2A'c - B')^2 - p^2D = 4A'(A'c^2 - B'c + C') \equiv 0 \pmod{p^2}.
\]
Ist $p = 2$, so ist $D'$ und folglich $B'$ gerade und mithin $B' - 2A'c$ immer durch 2 teilbar. Es ist dann
\[
\left(A'c - \frac{B'}{2}\right)^2 - D = A'(A'c^2 - B'c + C'),
\]
und da $D \equiv 0$ oder $\equiv 1 \pmod{4}$ ist, so ist die linke Seite hier durch 4 teilbar, wenn $A'c - \frac{B'}{2} \equiv D \pmod{2}$ angenommen wird.

Geht man umgekehrt von einer Gleichung
\weq{9}{A\omega^2 + B\omega + C = 0, \quad B^2 - 4AC = D}
der Diskriminante $D$ aus, nimmt $A$ relativ prim zu $p$ an und setzt $\omega' = p\omega$, so genügt $\omega'$ der primitiven Gleichung
\weq{10}{A\omega'^2 + Bp\omega' + Cp^2 = 0}
von der Diskriminante $D' = p^2D$, und wenn wir also $v = j(\omega')$
setzen, so ist $u = j(\omega)$ rational durch $v$ ausdrückbar. Daraus folgt also der Satz:

1. \gesperrt{Jede Klasseninvariante für die Diskriminante $D$ ist rational ausdrückbar durch eine Klasseninvariante für die Diskriminante $p^2D$.}

Um aber die Frage zu beantworten, wie viele Werte von $v$ zu demselben Werte von $u = j(\omega)$ führen, bemerken wir, daß die verschiedenen Werte von $v$, die dies leisten, Wurzeln der Gleichung (5) und (3) sein müssen, also in einer der Formen
\weq{11}{j(p\omega), \quad j\left(\frac{\omega + c}{p}\right)}
%CONT>>
